\documentclass[UTF8]{ctexart}
\title{基于自适应稀疏路由的长上下文大语言模型推理加速机制研究}
\begin{document}
\maketitle

\begin{abstract}
长上下文大语言模型在自回归解码阶段易出现键值缓存显存膨胀并拖慢服务响应。本文通过多尺度推理对比实验（$n = 48$），定量评估自适应稀疏路由（AdaptRoute）对长程检索任务的加速效果，并分析注意力熵阈值与分块粒度的影响规律。结果表明，在推荐配置下，AdaptRoute 使并发吞吐量提高 2.35 倍，单令牌解码延迟下降 68.4\%（$p < 0.01$），验证了低损动态缓存压缩的可行性。
\end{abstract}

\section{引言}
长文档理解与代码库级推理对大语言模型的上下文容量提出了更高要求~\cite{dao2022neurips,vaswani2017neurips}。传统全量注意力缓存占用显存过高，且易受带宽限制。依据层级注意力熵动态筛选关键键值块，可在低精度损失条件下将远距离证据令牌保留在高速缓存中（图~\ref{fig:route_cn}）。

\section{实验设计与分析方法}
实验采用开源 7B 与 70B 模型，在 128K 上下文与 BF16 精度下设置三组平行评测。使用硬件性能计数器分析压缩前后显存带宽利用率，并测定端到端解码延迟 $L_r$（$\text{ms token}^{-1}$）。

\begin{equation}
L_r = E_0 \max\!\left(0,\, \frac{\tau_b - \tau_c}{\tau_c}\right)^{1.5}
\end{equation}

\section{结果与讨论}
如表~\ref{tab:cn_res} 所示，两阶段自适应预算校准后，128K 场景下的相对困惑度偏移降至 4.6\%，多跳检索召回率由滑动窗口对照组的 29.5\% 提高至 87.6\%。消融监测显示，当保留比例维持在 6.0\% 以上时，深层检索头的长程注意力连贯性未见明显退化。
\end{document}
